\documentclass[10pt,a4paper]{amsart}
\usepackage[margin=19mm]{geometry}
\usepackage{amsmath,amssymb,mathtools}
\usepackage[scr=rsfs,cal=euler]{mathalfa}
\usepackage{xfrac}
\usepackage[unicode,hidelinks,bookmarks=false]{hyperref}
\newcommand{\ie}{i.e.\ }
\newcommand{\cf}{cf.\ }
\newcommand{\resp}{respectively\ }
\newcommand{\Subsection}{\subsection}
\providecommand{\nobreakdash}{\nobreak\hbox{-}}
\setlength{\emergencystretch}{2em}
\multlinegap0pt
\usepackage{amssymb}
\usepackage{mathtools}
\newtheorem{theorem}{Theorem}
\newtheorem{lemma}{Lemma}
\usepackage{cleveref}
\crefname{theorem}{Theorem}{Theorems}
\crefname{lemma}{Lemma}{Lemmas}
\newcommand{\Hyp}{\mathbb H}
\newcommand{\Mass}{\mathbf M}
\newcommand{\cur}[1]{\left[\!\left[#1\right]\!\right]}
\newcommand{\dist}{\operatorname{dist}}
\newcommand{\restr}{\mathbin{\llcorner}}
\newcommand{\supp}{\operatorname{supp}}
\title{Mass Bounds for Confined Area-Minimizing Minimal Surfaces}
\author{Xumin Jiang, Jiongduo Xie}
\date{Source: arXiv:2609.01324v2 (CC BY 4.0)}
\begin{document}
\maketitle

\noindent\emph{Source and licence.} Mass Bounds for Confined Area-Minimizing Minimal Surfaces. Version arXiv:2609.01324v2 (CC BY 4.0), distributed by arXiv under the Creative Commons Attribution 4.0 International licence, http://creativecommons.org/licenses/by/4.0/, as recorded in the arXiv licence metadata for this exact version and shown on the abstract page https://arxiv.org/abs/2609.01324v2. Full licence notice: \texttt{licenses/arxiv-2609.01324-CC-BY-4.0.txt}.\par\smallskip

\noindent\textit{Editorial scope.} Counted unit: the article's lemma lem:squashing (source0.tex lines 1285--1297). The article's complete original proof of it is delivered with it (source0.tex lines 1299--1454, one \texttt{proof} environment of 156 lines). No mathematics is written, repaired or re-derived: the statement and the proof are the article's own lines, in the article's own notation, and no retained line is substituted or re-flowed.  Retained uncounted as the source-local context the proof needs: lines 304--308; lines 317--342; lines 329--335; lines 481--484; lines 1245--1266; lines 1260--1263.  Nothing was omitted from any retained span, so the package carries no omission record.  No retained line contains a citation, so the wrapper carries no bibliography and no bibliography entry.  No retained line contains a figure environment, an image-inclusion command or drawing code, so no image is delivered and the package declares no asset dependency.  Editorial formatting only, in addition: an AMS \texttt{amsart} preamble that restates the article's own declarations and macros used by the retained lines, namely the environments \texttt{theorem}, \texttt{lemma} on separate counters, together with the standard packages those lines need.  The retained environments print in this document as Theorem 1, Lemma 1 in order of appearance, whereas the article prints the same items with the numbering of its own full document.  The complete native source of the article as distributed in the arXiv e-print for this exact version is bundled unchanged as \texttt{sources/209101/source0.tex}.
\begin{equation}\label{eq:vertical-plane-projection}
 L=\{z=0\}\cong\Hyp^n,
 \qquad
 \Pi_L(w,z,y)=(w,0,y).
\end{equation}

\begin{theorem}
\label{thm:local-mass}
There exist $\delta_*=\delta_*(n)>0$ and $C_*=C_*(n)<\infty$ with the
following property.  Let $T$ be an integer-rectifiable $n$-current with
locally finite mass in $\Hyp^{\mathfrak m}$, locally area minimizing, and let
$q\in\Hyp^{\mathfrak m}$.  Assume
\begin{equation}\label{eq:no-boundary-local}
 (\partial T)\restr B_2(q)=0.
\end{equation}
Suppose that, after an ambient isometry, $L$ and $\Pi_L$ have the form
\eqref{eq:vertical-plane-projection}.  Put $p=\Pi_L(q)$.  If, for some
integer $Q\ge1$ and $0<\delta\le\delta_*$,
\begin{align}
 \dist_{\Hyp}(q,L)&\le\delta,\label{eq:center-local}\\
 \supp T\cap B_2(q)&\subset
 \{\dist_{\Hyp}(\cdot,L)\le\delta\},\label{eq:tube-local}\\
 (\Pi_L)_\#\bigl(T\restr B_2(q)\bigr)\restr D_{3/2}^L(p)
 &=Q\,\cur{D_{3/2}^L(p)},\label{eq:projection-local}
\end{align}
then
\begin{equation}\label{eq:mass-local}
 \Mass_{\Hyp^{\mathfrak m}}\bigl(T\restr B_1(q)\bigr)\le C_*Q.
\end{equation}
The constants are independent of $k$, and no a priori upper bound for
$\Mass_{\Hyp^{\mathfrak m}}(T\restr B_2(q))$ is assumed.
\end{theorem}

\begin{align}
 \dist_{\Hyp}(q,L)&\le\delta,\\
 \supp T\cap B_2(q)&\subset
 \{\dist_{\Hyp}(\cdot,L)\le\delta\},\\
 (\Pi_L)_\#\bigl(T\restr B_2(q)\bigr)\restr D_{3/2}^L(p)
 &=Q\,\cur{D_{3/2}^L(p)},
\end{align}

\begin{equation}\label{eq:local-minimizer-definition}
 \Mass_U(T)\le \Mass_U(T+W),
 \qquad \Mass_U(S):=\|S\|(U).
\end{equation}

\begin{lemma}\label{lem:vertical-projection}
Let $x=(w,z,y)$ and let
$\rho(x)=\dist_{\Hyp}(x,L)$.  Then
\begin{equation}\label{eq:vertical-distance-formulas}
 \sinh\rho(x)=\frac{|z|}{y},
 \qquad
 d_{\Hyp}(x,\Pi_Lx)
 =2\operatorname{arsinh}\!\left(\frac{|z|}{2y}\right).
\end{equation}
The map $\Pi_L$ is $1$-Lipschitz.  Moreover, if $\rho(x)\le\delta\le1$,
then
\begin{equation}\label{eq:projection-displacement}
 d_{\Hyp}(x,\Pi_Lx)\le2\delta.
\end{equation}
The horizontal homotopy
\begin{equation}\label{eq:horizontal-homotopy}
 \mathcal H_\tau(w,z,y)=(w,(1-\tau)z,y),
 \qquad 0\le\tau\le1,
\end{equation}
has spatial differential of norm at most one and hyperbolic speed at most
$2\delta$ on the $\delta$-tube about $L$.
\end{lemma}

\begin{equation}
 \mathcal H_\tau(w,z,y)=(w,(1-\tau)z,y),
 \qquad 0\le\tau\le1,
\end{equation}

\begin{lemma}[Localized squashing with algebraic multiplicity]
\label{lem:squashing}
Under the assumptions of \cref{thm:local-mass}, set
\[
 F(r)=\Mass_{\Hyp^{\mathfrak m}}\bigl(T\restr B_r(q)\bigr),
 \qquad r\in[1/2,1].
\]
There is $c_s=c_s(n)<\infty$ such that, for $\delta$ sufficiently small,
\begin{equation}\label{eq:squashing-integrated}
 F(1/2)
 \le\exp\!\left(-\frac{1}{2c_s\delta}\right)F(1)+c_sQ.
\end{equation}
\end{lemma}

\begin{proof}
For almost every $r\in[1/2,1]$, slicing gives the integral cycle
\[
 Z_r=\partial\bigl(T\restr B_r(q)\bigr)
     =\langle T,d_q,r\rangle
\]
and
\begin{equation}\label{eq:slice-derivative}
 \Mass_{\Hyp}(Z_r)\le F'(r).
\end{equation}
Put $p=\Pi_L(q)$.  By \cref{lem:vertical-projection}, every point in the
support of $T\restr B_2(q)$ satisfies
\begin{equation}\label{eq:two-delta-displacements}
 d_{\Hyp}(x,\Pi_Lx)\le2\delta,
 \qquad
 d_{\Hyp}(q,p)\le2\delta.
\end{equation}
Set
\[
 s=r-4\delta.
\]
After decreasing $\delta_*$, one has $s\ge1/4$.

If $x\in\supp Z_r$ and $y=\Pi_Lx$, then
\begin{equation}\label{eq:projected-annulus}
 \bigl|d_L(y,p)-r\bigr|
 \le d_{\Hyp}(x,y)+d_{\Hyp}(q,p)
 \le4\delta.
\end{equation}
Let
\[
 S_r=(\Pi_L)_\#\bigl(T\restr B_r(q)\bigr).
\]
If $y\in D_s^L(p)$ and
$x\in\supp(T\restr B_2(q))$ satisfies $\Pi_Lx=y$, then
\[
 d_{\Hyp}(x,q)
 \le d_{\Hyp}(x,y)+d_L(y,p)+d_{\Hyp}(p,q)
 <2\delta+(r-4\delta)+2\delta=r.
\]
It follows that restriction before or after cutting down to $B_r(q)$
gives the same projected current over $D_s^L(p)$:
\[
 S_r\restr D_s^L(p)
 = (\Pi_L)_\#\bigl(T\restr B_2(q)\bigr)\restr D_s^L(p).
\]
Here the reverse inclusion is immediate from $B_r(q)\subset B_2(q)$, and
the preceding estimate proves the nontrivial inclusion of the relevant
preimages.  Since $s<3/2$, the projected-current identity
\eqref{eq:projection-local} therefore implies
\begin{equation}\label{eq:Sr-inner}
 S_r\restr D_s^L(p)=Q\,\cur{D_s^L(p)}.
\end{equation}
By locality of the boundary in the interior of $D_s^L(p)$,
$(\Pi_L)_\#Z_r=\partial S_r$ has no support there.  Moreover, $S_r$ is
supported in $D_{r+4\delta}^L(p)$.  Thus
\[
 C_r=S_r-Q\,\cur{D_s^L(p)}
\]
is supported in the closed annulus
$\overline{D_{r+4\delta}^L(p)}\setminus D_s^L(p)$ and satisfies
\begin{equation}\label{eq:boundary-Cr}
 \partial C_r=(\Pi_L)_\#Z_r-Q\,\cur{\partial D_s^L(p)}.
\end{equation}

Apply the horizontal homotopy \eqref{eq:horizontal-homotopy} to $Z_r$ and
set
\[
 A_r=\mathcal H_\#\bigl(\cur{[0,1]}\times Z_r\bigr).
\]
The homotopy formula gives
\begin{equation}\label{eq:Ar-boundary}
 \partial A_r=(\Pi_L)_\#Z_r-Z_r,
\end{equation}
and the product-current estimate together with
\cref{lem:vertical-projection} gives
\begin{equation}\label{eq:Ar-mass}
 \Mass_{\Hyp}(A_r)\le c_s\delta\,\Mass_{\Hyp}(Z_r).
\end{equation}

We next work inside $L$.  Write
$y=\exp_p(u\vartheta)$ in polar coordinates and, on the full annulus
$s\le u\le r+4\delta$ containing both $\supp C_r$ and
$\supp((\Pi_L)_\#Z_r)$, define
\begin{equation}\label{eq:radial-homotopy}
 \mathcal R_\tau\bigl(\exp_p(u\vartheta)\bigr)
 =\exp_p\bigl(((1-\tau)u+\tau s)\vartheta\bigr).
\end{equation}
Thus $\mathcal R_1$ is radial retraction to $\partial D_s^L(p)$.  Since
$C_r$ is an $n$-current in the $n$-manifold $L$ and
$\mathcal R_1$ has rank at most $n-1$,
$(\mathcal R_1)_\#C_r=0$.  Applying $\mathcal R_1$ to
\eqref{eq:boundary-Cr} yields
\begin{equation}\label{eq:radial-degree}
 (\mathcal R_1)_\#(\Pi_L)_\#Z_r
 =Q\,\cur{\partial D_s^L(p)}.
\end{equation}
Set
\[
 K_r=\mathcal R_\#\bigl(\cur{[0,1]}\times(\Pi_L)_\#Z_r\bigr).
\]
The homotopy formula and \eqref{eq:radial-degree} give
\begin{equation}\label{eq:Kr-boundary}
 \partial K_r=Q\,\cur{\partial D_s^L(p)}-(\Pi_L)_\#Z_r.
\end{equation}
The radial displacement in \eqref{eq:radial-homotopy} is at most
$8\delta$.  Its spatial differential has norm at most one: the radial factor
is $1-\tau$, while the angular factor is
\[
 \frac{\sinh((1-\tau)u+\tau s)}{\sinh u}\le1.
\]
Since $\Pi_L$ is $1$-Lipschitz,
\begin{equation}\label{eq:Kr-mass}
 \Mass_{\Hyp}(K_r)
 \le c_s\delta\,\Mass_{\Hyp}((\Pi_L)_\#Z_r)
 \le c_s\delta\,\Mass_{\Hyp}(Z_r).
\end{equation}

With compatible orientations, the current
\[
 R_r=-A_r-K_r+Q\,\cur{D_s^L(p)}
\]
satisfies $\partial R_r=Z_r$.  All currents in this construction are
supported in $B_{3/2}(q)$ when $\delta$ is sufficiently small.  Indeed,
the horizontal homotopy moves a point by at most $2\delta$, while every
point in the radial homotopy has distance at most $r+4\delta$ from $p$;
using $d_{\Hyp}(p,q)\le2\delta$ and $r\le1$ gives the assertion after fixing,
for example, $\delta_*<1/12$.  Since
$s\le1$,
\begin{equation}\label{eq:competitor-mass}
 \Mass_{\Hyp}(R_r)
 \le c_s\delta\,\Mass_{\Hyp}(Z_r)
     +Q\operatorname{Vol}_{\Hyp^n}(D_1^L(p))
 \le c_s\delta\,\Mass_{\Hyp}(Z_r)+c_sQ.
\end{equation}
The current
\[
 W_r=R_r-(T\restr B_r(q))
\]
is a cycle compactly supported in $B_{3/2}(q)$.  Since
$\|T\|(\partial B_r(q))=0$, applying
\eqref{eq:local-minimizer-definition} with this $W_r$ and then arguing as in
the Euclidean proof gives $F(r)\le\Mass(R_r)$.  Together with
\eqref{eq:slice-derivative} and
\eqref{eq:competitor-mass}, this implies
\begin{equation}\label{eq:differential-squashing}
 F(r)\le c_s\delta F'(r)+c_sQ
\end{equation}
for almost every $r\in[1/2,1]$.  Setting $G=F-c_sQ$, one obtains
\[
 G'(r)-\frac{1}{c_s\delta}G(r)\ge0.
\]
Hence $r\mapsto e^{-r/(c_s\delta)}G(r)$ is nondecreasing.  Evaluating at
$r=1/2$ and $r=1$ proves \eqref{eq:squashing-integrated}, after increasing
$c_s$ harmlessly.
\end{proof}

\end{document}
